%%%%%%%%%%%%%%%%%%%%%%%%%%%%%%%%%%%%%%%%%%%%%%%%%%%%%%%%%%%%%%
%%%%%%%%% MA-0361 Algebra Lineal I %%%%%%%%%%%%%%%%
%%%%%%%%%%%%%%%%%%%%%%%%%%%%%%%%%%%%%%%%%%%%%%%%%%%%%%%%%%%%%%
\newpage
\subsection*{MA-0361 Algebra Lineal I}
\addcontentsline{toc}{subsection}{MA-0361 Algebra Lineal I}

\begin{refsection}
\begin{table}[H]
\centering{
\begin{tabular}{|ll|}
\hline
%\multicolumn{2}{|c|}{\textbf{Nivel de virtualidad del curso:} Virtual}\\ \hline
\textbf{Año:} I  & \textbf{Requisitos:} Ninguno \\ 
\textbf{Ciclo:} II  & \textbf{Correquisitos:} Ninguno \\ 
\textbf{Tipo de curso:} Teórico & \textbf{Horas presenciales por semana:}   \\ 
\textbf{Créditos:} 3 & \textbf{Horas de trabajo independiente por semana:}  \\ \hline
\end{tabular}
}
\end{table}



\subsubsection*{Descripción del curso}

Este es el primero de una secuencia de dos cursos de álgebra lineal para la carrera de Matemática. Se enfoca en el desarrollo de habilidades operacionales y del pensamiento intuitivo y geométrico, con base en la exploración de una serie de conceptos del álgebra lineal. Adicionalmente, se introduce al estudiantado a un mayor nivel de abstracción y formalismo mediante la demostración de resultados elementales. 

El desarrollo de los temas inicia con conceptos concretos como sistemas de ecuaciones lineales, vectores, matrices y determinantes, y en la segunda parte del curso se les da un fundamento teórico a estos por medio de las nociones de espacios vectoriales y transformaciones lineales.

El curso brinda algunas herramientas básicas que se utilizan en el estudio del cálculo diferencial e integral en varias variables en el curso MA-0351 Introducción al cálculo en varias variables, así como en el resto de la carrera.

\subsubsection*{Objetivos}

\textbf{General:} Aplicar las nociones, intuiciones y conceptos básicos del álgebra lineal para la resolución de problemas. \\

\textbf{Específicos:} Al finalizar el curso se espera que la persona estudiante sea capaz de:

\begin{enumerate}
\item Resolver sistemas de ecuaciones lineales.
\item Realizar cálculos con los distintos objetos del álgebra lineal.
\item Relacionar los conceptos de solubilidad de sistemas, invertibilidad de matrices, determinantes y otras condiciones equivalentes.
\item Aplicar los conceptos básicos del álgebra lineal en la demostración de resultados y la resolución de problemas teóricos y prácticos.
\item Interpretar gráficamente los conceptos básicos del álgebra lineal.
\item Experimentar con los distintos objetos del álgebra lineal a través de software computacional. (SH03)
\item Indagar en la literatura sobre temas que permitan complementar su formación.
\item Comunicar ideas matemáticas de manera clara y efectiva en forma oral y escrita.
\end{enumerate}

\subsubsection*{Contenidos}

\begin{enumerate}
\item \textbf{Sistemas de ecuaciones lineales (2 semanas):} 
\begin{enumerate}
    \item Solución de sistemas de dos ecuaciones con dos incógnitas y su interpretación geométrica.
    \item Sistemas de ecuaciones $n \times m$ y eliminación gaussiana.
    \item Interpretación geométrica de la solución de un sistema de tres ecuaciones con tres incógnitas.
    \item Sistemas homogéneos de ecuaciones.
\end{enumerate}

\item \textbf{Vectores y matrices (3 semanas):} 
\begin{enumerate}
    \item Definiciones generales. Operaciones con matrices y vectores: suma, resta y multiplicación de matrices, producto escalar de vectores, producto de escalares por matrices.
    \item Matrices y sistemas de ecuaciones lineales.
    \item Inversa de una matriz cuadrada, matriz transpuesta, matrices elementales y matrices inversas.
    \item Factorizaciones LU de una matriz.
\end{enumerate}

\item \textbf{Determinantes (2 semanas):} 
\begin{enumerate}
    \item Definiciones.
    \item Interpretación geométrica de un determinante $2 \times 2$.
    \item Propiedades de los determinantes.
    \item Determinantes y matrices inversas.
    \item Regla de Cramer.
    \item Permutaciones y determinantes.
\end{enumerate}

\item \textbf{Vectores en $\R^2$ y $\R^3$ (1 semanas):} 
\begin{enumerate}
    \item Vectores en $\R^2$ y $\R^3$. magnitud de vectores, ángulo entre vectores, proyecciones ortogonales.
    \item El producto cruz de dos vectores.
    \item Rectas y planos en el espacio.
\end{enumerate}

\item \textbf{Espacios vectoriales (4 semanas):} 
\begin{enumerate}
    \item Definiciones y propiedades básicas.
    \item Subespacios vectoriales.
    \item Combinaciones lineales y susbespacios generados.
    \item Independencia lineal.
    \item Bases y dimensión.
    \item Cambio de bases.
    \item Rango, nulidad, espacios de filas y de columnas.
\end{enumerate}

\item \textbf{Transformaciones lineales (3 semanas):} 
\begin{enumerate}
    \item Definición y ejemplos. Interpretación geométrica en ejemplos concretos.
    \item Propiedades de las transformaciones lineales. 
    \item Imagen y núcleo de una transformación lineal.
    \item Representación matricial de una transformación lineal, cambio de bases y similaridad.
    \item Isomorfismos: inyectividad y sobreyectividad de transformaciones lineales, isomorfismos de espacios vectoriales.
\end{enumerate}

\end{enumerate}

\subsubsection*{Metodología}

\noindent \textbf{Lineamientos metodológicos:} La persona docente a cargo de este curso debe respetar las siguientes pautas:
\begin{enumerate}
    \item La presentación de los contenidos debe priorizar la intuición y el desarrollo de habilidades operacionales por encima del formalismo. 
    \item Se deben realizar las demostraciones que sean de caracter operacional y también las que involucren argumentos lógicos sencillos. Por ejemplo, resultados sobre independencia y dependencia lineal, subespacios, transformaciones lineales. Por el contrario, se deben evitar demostraciones de alta complejidad o abstracción, como por ejemplo la demostración de que todo espacio vectorial tiene una base.
    %\item Como existe codependencia entre la temática y objetivos de este curso y la del curso MA-02XX Introducción al Cálculo en Varias Variables, debe respetarse estrictamente el cronograma establecido en este documento.
\end{enumerate}

\textbf{Sugerencias metodológicas:} Adicionalmente, se tienen las siguientes recomendaciones para la persona docente a cargo de este curso:
\begin{enumerate}
    \item Utilizar el recurso tecnológico para reforzar distintos conceptos estudiados en el curso por medio de la visualización, cálculo y experimentación. En particular, se podría asignar alguna tarea o proyecto que requiera utilizar un sistema algebraico computacional (CAS) para realizar algunas de estas tareas.
    \item Aprovechar momentos adecuados del curso para motivar el estudio por medio de reseñas acerca del desarrollo histórico de la matemática y su importancia en aplicaciones. Por ejemplo, el libro de Grossman \cite{Gro19} y el de Poole \cite{Poo17} contienen reseñas y semblanzas históricas, las cuales podrían ser asignadas como lecturas a las personas estudiantes o utilizadas por la persona docente para motivar la introducción de algún tema del curso.
    \item En la evaluación, se sugiere utilizar estrategias que promuevan el trabajo constante de parte de las personas estudiantes a lo largo del ciclo lectivo. Por ejemplo, esto se puede hacer por medio de tareas o quizzes.
    \item Se recomienda incorporar durante el proceso de aprendizaje, ejercicios como los que aparecen en las subsecciones denominadas \textit{Exploración} en el libro de Poole \cite{Poo17}. También, se podrían asignar ensayos como los planteados en las secciones denominadas \textit{Proyecto de ensayo} del mismo libro.
    \item Dado que hoy en día la mayoría del trabajo en matemática, tanto a nivel académico como profesional, es de naturaleza colaborativa, se sugiere a la persona docente implementar estrategias que promuevan el trabajo en equipo.
\end{enumerate}




\nocite{Apo84}
\nocite{Apo85} 
\nocite{Gro19}
\nocite{MM18}
\nocite{HK71}
\nocite{Cur93}
\nocite{Lan87}
\nocite{Tre17}



\printbibliography[heading=subbibliography,section=\the\value{refsection}]
\end{refsection}
